% results.tex (v6): Results, six subsections, about 2300 words. All numbers are macros from numbers.tex.
\section{Results}\label{sec:results}

\bmhead{A cell-year test bed with an explicit time definition}
The data set holds \nncyAll{} cell-years in \nncellAll{} 1 km cells (Fig.~1b--d; Table~1). The three regions differ in how their observations are spread over years (Fig.~1e). In Alaska \ntopyrshareAk{}\% of the cell-years fall in \ntopyrAk{}, and \nnmultiAk{} cells were observed in four or more years; the Lena Delta cell-years are spread over \nnyrLe{} years with \nnmultiLe{} multi-year cells; Canada has \nncyCa{} cell-years concentrated in \ntopyrCa{} with \nnmultiCa{} multi-year cells. Within one season, thaw depth at the Lena Delta sites grows in proportion to the square root of the degree-days accumulated to the measurement day (Fig.~2a), which is the Stefan relation and the reason the forcing is accumulated to the observation day rather than over the full year. Across years at the same cell, the relation is weaker: cells observed in several years show year-to-year changes that only partly follow the change in degree-days (Fig.~2b).

\bmhead{Year-specific forcing helps where observations span many years}
Replacing climatological forcing by the forcing of the observation year changed the spatial-block RMSE of the recalibrated Stefan model by \ndforcLePo{} cm in the Lena Delta (Fig.~3a); the source-coefficient Stefan and Kudryavtsev models changed by \ndforcLePz{} and \ndforcLeKu{} cm. In Alaska the change was \ndforcAkPo{} cm and in Canada \ndforcCaPo{} cm, both with confidence intervals that include zero. The Alaska result is not evidence against year-specific forcing: with half of the labels in one year, the climatological and yearly forcings differ little for most training rows, and the spatial-block split cannot separate the year effect from the site effect. The learned models gained less from year-specific forcing than the physical models (direct learning \ndforcLeDz{}, physics plus residual \ndforcLeRo{} cm in the Lena Delta), because the covariates already carry part of the climate signal.

\bmhead{The validation split sets the ranking of methods}
With a random split by cell, direct learning had the lowest RMSE in all three regions (Fig.~3b): \nrmseAkrandDz{} cm in Alaska, \nrmseLerandDz{} cm in the Lena Delta and \nrmseCarandDz{} cm in Canada. Moving to the spatial-block split raised these values by \noptimAkDz{}, \noptimLeDz{} and \noptimCaDz{} cm, because the random split places other years of the same cell, and neighbouring cells of the same site, in the training set. The physical models changed little between splits (recalibrated Stefan \nrmseAkrandPo{} versus \nrmseAkblkPo{} cm in Alaska). Under block-by-year hold-out, in which both the block and the year of the scored cell-years are unseen, the advantage of the combined model over Stefan shrank to \ndrefLeblkyrRo{} cm in the Lena Delta and \ndrefCablkyrRo{} cm in Canada. Under region hold-out, in which a region is predicted from the other two, direct learning was no better than recalibrated Stefan in Alaska (\ndrefAkregDz{} cm) and far worse in the Lena Delta (\ndrefLeregDz{} cm; RMSE \nrmseLeregDz{} cm), whereas the physics-plus-residual model stayed within \ndrefLeregRo{} cm of Stefan in the Lena Delta and at \nrmseAkregRo{} cm in Alaska. The physical anchor is what keeps the combined model usable outside its training region.

\bmhead{Residual learning improves on Stefan under identifiable conditions}
On the spatial-block split with the full feature set, the physics-plus-residual model lowered RMSE relative to recalibrated Stefan by \ndrefLeblkRo{} cm in the Lena Delta and \ndrefCablkRo{} cm in Canada, with confidence intervals that exclude zero, and by \ndrefAkblkRo{} cm in Alaska, which is not distinguishable from zero (Fig.~4a). Direct learning gave \ndrefLeblkDz{} (Lena Delta), \ndrefCablkDz{} (Canada) and \ndrefAkblkDz{} cm (Alaska), physics-input learning \ndrefLeblkFo{}, \ndrefCablkFo{} and \ndrefAkblkFo{} cm, and physics pseudo-labels \ndrefLeblkDo{}, \ndrefCablkDo{} and \ndrefAkblkDo{} cm. Pseudo-labels therefore reproduced the information of the physical model without adding to it.

The gain depends on the covariates. With climate variables only, the residual model was no better than Stefan in any region and worse in Alaska (\ndfeatAkFclim{} cm); adding terrain gave \ndfeatAkFterr{} cm and adding soil \ndfeatAkFsoil{} cm in Alaska, with \ndfeatLeFsoil{} cm in the Lena Delta and \ndfeatCaFsoil{} cm in Canada (Fig.~4a). MODIS vegetation and snow-off date, high-resolution soil and land-cover classes and lagged forcing changed the gain by less than 1 cm (Alaska \ndfeatAkFhires{} cm with the high-resolution set). The gain also depends on the thermal regime: within regional terciles of thawing degree-days, the residual model improved on Stefan most in the warmer terciles (Lena Delta warm tercile \ndtddLewarm{} cm, Canada mid tercile \ndtddCamid{} cm) and least in the cold terciles (Lena Delta \ndtddLecold{}, Canada \ndtddCacold{} cm), while in Alaska the change was within 1 cm in every tercile (Fig.~4b). Replacing CatBoost by random forest or by three neural networks changed the gain by less than the differences between regions and did not change its sign in any region (Fig.~4c); in Alaska the neural networks were slightly worse than Stefan (multilayer perceptron \ndlrnAkmlpRo{} cm), and in Canada direct learning with FT-Transformer (\ndlrnCafttDz{} cm) was within the interval of CatBoost direct learning (\ndlrnCacatboostDz{} cm).

\bmhead{Same-cell, same-year comparison with published products}
On the common cell-years with values in every product (Fig.~4d), the physics-plus-residual model had an RMSE of \nccLeRoall{} cm in the Lena Delta against \nccLeWeiall{} cm for Wei et al., \nccLeLiuall{} cm for Liu et al. and \nccLeCcifourall{} cm for ESA CCI v4; the paired difference against Wei et al. was \ndprodLeweiRomask{} cm. In Alaska the comparison is affected by leakage: Wei et al. and Liu et al. were trained on Alaska observations that overlap ours. On all common cells our model gave \nccAkRoall{} cm against \nccAkWeiall{} cm for Wei et al. and \nccAkLiuall{} cm for Liu et al. After removing cells within 25 km of the two products' training sites, our model gave \nccAkRomask{} cm, Wei et al. \nccAkWeimask{} cm and Liu et al. \nccAkLiumask{} cm; the paired differences were \ndprodAkweiRomask{} cm against Wei et al. and \ndprodAkliuRomask{} cm against Liu et al. In Canada direct learning gave \nccCaDzall{} cm and the residual model \nccCaRoall{} cm against \nccCaWeiall{} cm for Wei et al. (paired difference \ndprodCaweiRomask{} cm). ESA CCI v5 and Aalto et al. had RMSEs above 45 cm in all regions because their values are far from the observed range. Calibration curves (Fig.~4e) show that all learned models compress the observed range: deciles above 75 cm are under-predicted in all regions.

\bmhead{Year-wise behaviour}
Under year hold-out the correlation between predicted and observed interannual anomalies was \nranomAkyrRo{} (Alaska), \nranomLeyrRo{} (Lena Delta) and \nranomCayrRo{} (Canada) for the residual model and \nranomAkyrPo{}, \nranomLeyrPo{} and \nranomCayrPo{} for recalibrated Stefan (Fig.~3c); the residual learner did not add interannual skill to the physical anchor. For cells observed in several years, the correlation between observed and predicted year-to-year changes was \nrdeltaAkRo{} for the residual model in Alaska (\nnpairsAk{} pairs) against \nrdeltaAkWei{} for Wei et al. and \nrdeltaAkLiu{} for Liu et al.; in the Lena Delta (\nnpairsLe{} pairs) the products reached \nrdeltaLeWei{} (Wei et al.) and \nrdeltaLeLiu{} (Liu et al.) against \nrdeltaLeRo{} for the residual model, and in Canada (\nnpairsCa{} pairs) no method exceeded \nrdeltaCaLiu{} (Fig.~3d). Year-wise bias of the residual model under the spatial-block split ranged from \nybiasminAk{} to \nybiasmaxAk{} cm in Alaska, \nybiasminLe{} to \nybiasmaxLe{} cm in the Lena Delta and \nybiasminCa{} to \nybiasmaxCa{} cm in Canada (Fig.~3e); with full-year instead of observation-day forcing the Alaska range was \nybiasminAkYr{} to \nybiasmaxAkYr{} cm (standard deviation across years \nybiassdAkYr{} against \nybiassdAk{} cm), which is the direct effect of matching the forcing to the observation day. The largest remaining biases coincide with changes in campaign composition rather than with climate: the Lena Delta 2020 campaign sampled channel sites, and the Canada 2023--2024 campaigns sampled a burned catchment. Within the same cell-year, measurements made in September are deeper than those made in August by about 10 cm, while probe and radar measurements did not differ (Supplementary Information). At the longest-observed sites (Fig.~5) the yearly predictions stay within the observed interannual range, the predicted year-to-year spread is smaller than the observed spread at most sites (Fig.~5g), and the published products either miss the site level or show no interannual variation.

\bmhead{Year-labelled 1 km maps}
Fitted on all cell-years of a region, the physics-plus-residual model was applied to the 1 km grid with the ERA5-Land forcing of each year 2011--2024 (Fig.~6). Each map carries the year of its forcing: in Alaska the 2014 map is shallower than the 2011--2024 mean over the North Slope and the 2022 map is deeper over the interior. Thaw depth follows terrain within the 20 km windows, with deeper values on south-facing slopes and valley floors. The area of applicability covers \naoaAk{}\% of the Alaska grid and \naoaLe{}\% of the Lena Delta grid; cells outside are shown faded.
